% !TEX root = ../Report.tex
% ======================================================================
\chapter{Requirements and Project Specification}
\label{ch:requirements}
% ======================================================================

This chapter describes the requirements that guided GiTrip's development, how
they were derived, and how they map to success criteria and evaluation. It also
documents scope boundaries, including the explicit decision to exclude Artificial Intelligence (AI)-powered
itinerary generation from the Minimum Viable Product (MVP).

\section{Requirements Elicitation}
Requirements were formed using two complementary inputs:
\begin{itemize}
  \item \textbf{Background research} (from Chapter~\ref{ch:background}) into competitor tools (maps, generic collaboration, and trip-planning products) and foundational techniques (version control, scheduling heuristics, etc.).
  \item \textbf{A preliminary survey} (as discussed in Section~\ref{sec:motivation} and Section~\ref{sec:overview}) to understand current trip-planning workflows, perceived pain points, and interest in proposed features.
\end{itemize}

\section{MoSCoW Objectives}
The Project Specification defined MoSCoW objectives. The following is the summary of Must-have (M1--M5) and Should-have (S1--S2) objectives, which formed the MVP definition:

\subsection*{Must-haves (M)}
These five objectives address the core gap identified in Chapter~\ref{ch:background}: no existing tool combines version control, constraint-aware scheduling, visual editing, and structured merging in a single workflow.
\begin{enumerate}[label=\textbf{M\arabic*}]
  \item \textbf{Trip repositories with history:} commits, branches, and snapshot history; visualisation of version history. This directly addresses the lack of change management in existing trip planners (Section~\ref{sec:competitors}).
  \item \textbf{Optimised automatic scheduler:} produce a feasible multi-day plan from places, durations, ordering constraints, desired windows, active hours, breaks, target day ranges, and preferences. The preliminary survey found 60\% of respondents reported planning as often or always stressful, motivating automation of the most tedious step.
  \item \textbf{Routing data integration:} retrieve travel-time data (e.g., OpenRouteService and Google Directions) and feed into scheduling. Without real travel durations, the scheduler cannot produce feasible plans.
  \item \textbf{Interactive visualisation:} timeline and map with editing and committing. The survey's feature ranking (Q5) placed timeline/map visualisation as the most valued capability.
  \item \textbf{Merge and conflicts:} structured merge for plan JSON with a conflict resolver screen. This is the core mechanism that enables collaborative editing without accidental overwrites.
\end{enumerate}

\subsection*{Should-haves (S)}
Should-haves improve the quality of generated plans and the branching workflow but are not essential for a demonstrable MVP:
\begin{enumerate}[label=\textbf{S\arabic*}]
  \item \textbf{Opening-hours validation:} automatic nudge to nearest feasible slot.
  \item \textbf{User preferences:} compact vs.\ sparse scheduling; morning/midday/night focus; recompute travel by mode (driving/walking/cycling/transit).
  \item \textbf{Branch productivity:} create branch from Web UI, compare branches, and fork plan from a public repository.
\end{enumerate}

Several Should-have and Could-have features were also added as direct responses to preliminary survey feedback (Q7 and Q8 in Appendix~\ref{appx:prelim_survey}). For example, weather integration, travel disruption alerts, packing checklists, and calendar export were requested by respondents.

\subsection*{Could-haves (C)}
Could-haves would enhance the user experience but were deprioritised to protect the MVP schedule:
\begin{enumerate}[label=\textbf{C\arabic*}]
  \item \textbf{Command-line interface (CLI):} Git-style terminal commands alongside the Web UI.
  \item \textbf{Public trip repository:} public gallery with attribution and one-click fork.
  \item \textbf{Real-time tracking:} live GPS guidance to detect running late and suggest reflow.
  \item \textbf{Busy-time consideration:} account for crowd levels in the auto-scheduler.
  \item \textbf{Progressive Web App (PWA):} offline read-only access and queued edits.
\end{enumerate}

\subsection*{Won't-haves (W)}
These were excluded to keep the project focused on deterministic, inspectable scheduling and version control semantics. Each would introduce substantial scope expansion without strengthening the core contribution:
\begin{enumerate}[label=\textbf{W\arabic*}]
  \item \textbf{Real-time collaborative editing:} simultaneous live editing via Conflict-free Replicated Data Type (CRDT) or Operational Transformation (OT).
  \item \textbf{Booking/reservation integration:} interfacing with hotel, restaurant, or ticket systems.
  \item \textbf{Budget tracking and expense splitting:} cost estimates, shared expenses, and currency conversion.
  \item \textbf{Social features:} chat, reviews, and ratings within the app.
  \item \textbf{Mobile native apps:} iOS/Android apps beyond PWA.
  \item \textbf{AI-powered itinerary generation:} natural-language trip generation (e.g., ``plan a romantic weekend in Paris'').
  \item \textbf{Learning and smart assistance:} preference learning to bias the auto-scheduler or recommend places from past trips.
\end{enumerate}

\section{Scope Boundary: Why AI Was Excluded}
\label{sec:ai_scope}
The second assessor noted in the Progress Report feedback:
\begin{quote}
``Since this is one of the very few projects I have seen without explicit use of AI, I would suggest student to investigate how some Machine Learning techniques could be exploited to make the project even stronger.''
\end{quote}
This suggestion is acknowledged. However, AI-powered itinerary generation was deliberately excluded from the Project Specification (W6) at an early stage for four practical reasons:
\begin{enumerate}
    \item determinism and explainability (the scheduler must justify feasibility decisions with rule-based logic)
    \item evaluation clarity (success criteria require objective
    measurement of constraint satisfaction and merge usability)
    \item deployment simplicity (avoiding model hosting and preference profiling), and 
    \item risk control (AI would broaden the problem into recommendation and
    natural-language understanding, exceeding the MVP schedule).
\end{enumerate}

That said, AI remains an interesting direction for future work, and potential extensions such as preference learning, dwell-time prediction, and context-aware recommendations are discussed in Chapter~\ref{ch:conclusions}.

\section{Success Criteria}
The Project Specification defined three success criteria, each tied to a Must-have objective. Thresholds were set as pragmatic targets that balance ambition with the constraints of performing UAT on the Evaluation Day (limited participant pool and 40-minute demonstration slot).

\begin{enumerate}[label=\textbf{SC\arabic*}]
  \item \textbf{Merge usability (M5):} create two alternative plans, merge, and resolve at least one timing conflict in $<5$ minutes on average by 80\% of 10--30 test participants.
  \item \textbf{Scheduling feasibility (M2):} automatic scheduler produces a non-overlapping sequence of activities for $\ge 6$ stops across 1--2 days with non-zero, realistic travel gaps.
  \item \textbf{Usability and accessibility (M4):} no critical issues with accessibility and follows WCAG 2.2.
\end{enumerate}

\section{Non-functional Requirements}
In addition to the functional objectives above, the following non-functional requirements were maintained throughout development given external Application Programming Interface (API) dependencies and deployment needs:

\begin{enumerate}[label=\textbf{NF\arabic*}]
  \item \textbf{Reliability:} graceful fallback when routing APIs fail; avoid
        corrupting repository state.
  \item \textbf{Performance:} responsive UI; reasonable planning time for
        typical stop counts (single-digit to low-teens).
  \item \textbf{Security and privacy:} session cookies; rate limiting; input
        validation; avoid collecting unnecessary personal data; General Data Protection Regulation (GDPR)~\cite{ref_gdpr} awareness.
  \item \textbf{Deployability:} simple deployment model with persistent storage
        (SQLite) and environment variables for API keys.
\end{enumerate}
